\section{Console 输入层：\texttt{console.c}}

\codefile{src/kernel/core/console.c}
\begin{lstlisting}[style=cstyle]
int console_try_getc(char* out) {
    // 优先查键盘 IRQ 缓冲区（非阻塞）
    if (keyboard_try_getc(out)) return 1;

    // 其次查串口（轮询，非阻塞）
    if (serial_try_getc(out)) {
        if (*out == '\r') *out = '\n';  // 串口终端发 CR，归一化为 LF
        return 1;
    }
    return 0;  // 无输入
}

char console_getc(void) {
    for (;;) {
        char c;
        if (console_try_getc(&c)) return c;
        __asm__ volatile("pause");  // 让出 CPU 一个周期
    }
}
\end{lstlisting}

\begin{whybox}
为什么 \texttt{console\_getc} 用 \texttt{pause} 而非 \texttt{hlt}？

\texttt{hlt} 会让 CPU 完全停止直到下一个中断。但 \texttt{serial\_try\_getc} 是
轮询模式（不依赖中断），\texttt{hlt} 后如果没有键盘中断就永远不会被唤醒。
\texttt{pause} 只是一个短暂提示（用于 spin-wait），CPU 仍然继续执行循环。

更好的方案是在每次 \texttt{hlt} 前 check 一次串口，但当前简单实现已够用。
\end{whybox}

% ============================================================================

